% =====================================================================
%  WORKING REPORT
%  Consolidated research notes on one-electron diagnostics for
%  machine-learned exchange-correlation functionals
%
%  This is a *working document*, not a submission. It is meant to be
%  read, edited, and re-issued as the research progresses.
% =====================================================================
\documentclass[11pt]{report}

\usepackage[margin=1in]{geometry}
\usepackage{amsmath,amssymb,amsthm}
\usepackage{graphicx}
\usepackage{booktabs}
\usepackage{longtable}
\usepackage{hyperref}
\usepackage[T1]{fontenc}
\usepackage{lmodern}
\usepackage{enumitem}
\usepackage{xcolor}
\usepackage{framed}

\hypersetup{colorlinks=true,linkcolor=blue,citecolor=blue,urlcolor=blue}

% ---- theorem environments (used sparingly; most content is plain prose) ----
\newtheorem{theorem}{Theorem}
\newtheorem{lemma}{Lemma}
\newtheorem{proposition}{Proposition}
\newtheorem{corollary}{Corollary}
\theoremstyle{definition}
\newtheorem{definition}{Definition}
\theoremstyle{remark}
\newtheorem{remark}{Remark}

% ---- custom commands for this document ----
\newcommand{\Ha}{\text{Ha}}
\newcommand{\mHa}{\text{mHa}}
\newcommand{\ang}{\text{\AA}}
\newcommand{\SIE}{\mathrm{SIE}}
\newcommand{\Exc}{E_{xc}}
\newcommand{\EHF}{E_{\mathrm{HF}}}
\newcommand{\Eexact}{E_{\mathrm{exact}}}
\newcommand{\ECD}{E_{\mathrm{CIDER}}}
\newcommand{\conv}{\mathrm{conv}}
\newcommand{\stable}{\mathrm{stable}}

% ---- confidence tags ----
\newcommand{\verified}[1]{\textbf{[V]}~#1}      % verified numerically, reproducible
\newcommand{\tentative}[1]{\textbf{[T]}~#1}      % plausible, not yet fully checked
\newcommand{\withdrawn}[1]{\textbf{[W]}~#1}      % previously claimed, now doubted or retracted
\newcommand{\todo}[1]{\textcolor{red}{\textbf{[TODO:}~#1\textbf{]}}}

\title{%
  \textbf{One-Electron Diagnostics for Machine-Learned\\ Exchange--Correlation Functionals}\\[0.5em]
  \large A Working Report: Theory, Calculations, Claims, and Open Problems%
}
\author{Dalton Omondi}
\date{Working draft --- \today}

\begin{document}
\maketitle

\begin{center}
\fbox{\parbox{0.9\textwidth}{
\textbf{Status of this document.}
This is a \emph{working report}, not a manuscript submission.
It consolidates five earlier drafts into a single narrative, states every
claim with an explicit confidence tag, and lists every known problem that
must be resolved before any part of this work is submitted for publication.
The tags used throughout are: \verified{verified numerically and
reproducible}; \tentative{plausible but not yet fully checked};
\withdrawn{previously claimed but now doubted or retracted}.
Passages marked \todo{} are unresolved.
}}
\end{center}

\tableofcontents
\clearpage

% =====================================================================
\chapter{Purpose and scope}
% =====================================================================

\section{What this document is}

This report consolidates five earlier working drafts into one narrative.
The five drafts were:

\begin{enumerate}[label=\textbf{D\arabic*.}]
  \item Hartree--Fock exactness for one-electron systems and its use as a
        basis-consistent reference for self-interaction diagnostics.
  \item A same-basis reproduction of self-interaction and delocalization
        errors in PBE, r2SCAN, and B3LYP.
  \item A symmetry-breaking bifurcation in the CIDER26XCCHEM machine-learned
        exchange--correlation functional on the H$_2^+$ dissociation curve.
  \item Formulation and refutation of a ``one-electron-limit patch''
        hypothesis for machine-learned functionals.
  \item A categorical comparison of Skala and CIDER as convergent
        constraint-encoding designs.
\end{enumerate}

The drafts were written in the order above but are not independent. D1
establishes the reference used by D2 and D3. D3 motivates D4. D5 tries to
abstract D2 and D3 into a design-principle comparison. Presenting them
separately as five papers creates circular citations (each cites the others
as ``companion manuscript''), cross-paper numerical inconsistencies, and a
false impression of five independent contributions when there is at most
one.

The honest summary, stated at the outset so that the rest of this report can
be read with the right expectations, is:

\begin{quote}
\textbf{The scientific content of the five drafts reduces to one
potentially publishable finding --- the symmetry-breaking bifurcation in
CIDER26XCCHEM on stretched H$_2^+$ --- supported by two small methodological
points (the same-basis reference, and the non-totality of the
constraint-residual order), one reproduction study, and two trivialities
(the patch refutation, and the categorical framing of the constraint
comparison). Whether the bifurcation survives a proper SCF stability
analysis is the single question that determines whether this work is a
paper or a methods note.}
\end{quote}

Everything below is written to serve that determination.

\section{What this document is not}

It is not a submission draft. It does not attempt to satisfy journal style,
length limits, or referee expectations. It includes working notes,
unresolved questions, and explicit doubts that would not appear in a
published paper.

It is also not a review of the field. It assumes no prior knowledge of
density functional theory (DFT) or of machine-learned functionals, and
explains both from first principles, but it does not attempt to survey the
literature beyond the specific papers cited.

\section{How to use it}

\begin{itemize}
  \item \textbf{To understand the project}: read Chapters~\ref{ch:background}
        through \ref{ch:results} in order.
  \item \textbf{To decide what to do next}: read Chapter~\ref{ch:issues}
        (known problems) and Chapter~\ref{ch:future} (next actions).
  \item \textbf{To communicate the work}: the claims in
        Chapter~\ref{ch:claims} are the load-bearing statements; every
        other chapter is scaffolding for them.
  \item \textbf{To write a paper}: Section~\ref{sec:paper-plan} at the end
        gives the recommended paper structure if the bifurcation survives
        verification.
\end{itemize}

% =====================================================================
\chapter{Background}
\label{ch:background}
% =====================================================================

This chapter assumes no prior knowledge of DFT. A reader already fluent in
the subject can skim and skip to Chapter~\ref{ch:theory}.

\section{Density functional theory in one page}

The central object of quantum chemistry is the many-electron Schr\"odinger
equation,
\begin{equation}
\hat H \Psi = E \Psi,
\qquad
\hat H = -\tfrac12 \sum_i \nabla_i^2
        - \sum_{i,A} \frac{Z_A}{|\mathbf r_i - \mathbf R_A|}
        + \sum_{i<j} \frac{1}{|\mathbf r_i - \mathbf r_j|},
\label{eq:many-body-h}
\end{equation}
where $i,j$ label electrons, $A$ labels nuclei, $Z_A$ is a nuclear charge,
and atomic units are used ($\hbar = m_e = e = 1$). The wavefunction
$\Psi(\mathbf r_1, \dots, \mathbf r_N)$ depends on $3N$ spatial coordinates,
so direct solution is intractable beyond a handful of electrons.

DFT replaces the wavefunction with the electron density
\begin{equation}
n(\mathbf r) = N \int |\Psi(\mathbf r, \mathbf r_2, \dots, \mathbf r_N)|^2
               \, d\mathbf r_2 \cdots d\mathbf r_N,
\end{equation}
a function of only three coordinates. The Hohenberg--Kohn theorem (1964)
states that the ground-state density uniquely determines the external
potential, and therefore the Hamiltonian and all ground-state properties.
Kohn and Sham (1965) then introduced a fictitious non-interacting system
with the same density as the real system. Its wavefunction is a single
Slater determinant of orbitals $\varphi_i$, and the density is
$n(\mathbf r) = \sum_i |\varphi_i(\mathbf r)|^2$.

The Kohn--Sham total energy is written as
\begin{equation}
E[n] = T_s[n] + V_{\rm ext}[n] + J[n] + \Exc[n],
\label{eq:ks-energy}
\end{equation}
where
\begin{align}
T_s[n] &= -\tfrac12 \sum_i \int \varphi_i^*(\mathbf r) \nabla^2
          \varphi_i(\mathbf r)\, d\mathbf r
          && \text{(non-interacting kinetic energy)},\\
V_{\rm ext}[n] &= \int v_{\rm ext}(\mathbf r)\, n(\mathbf r)\, d\mathbf r
          && \text{(nuclear attraction)},\\
J[n] &= \tfrac12 \iint \frac{n(\mathbf r_1)\, n(\mathbf r_2)}
                          {|\mathbf r_1 - \mathbf r_2|}
        \, d\mathbf r_1 \, d\mathbf r_2
          && \text{(classical Coulomb self-repulsion)}.
\end{align}
The exchange--correlation functional $\Exc[n]$ is defined by
Eq.~\eqref{eq:ks-energy} as \emph{whatever remains} once the first three
terms are subtracted from the exact total energy. It is not known in closed
form and must be approximated.

\section{The exchange--correlation functional and its exact constraints}

Because $\Exc$ is defined by subtraction, it inherits a set of exact
mathematical properties from the true ground-state energy. These are the
\emph{exact constraints} on $\Exc$. Well-known examples include:

\begin{itemize}
  \item \textbf{Uniform coordinate scaling.}
        For $\lambda > 0$, define $n_\lambda(\mathbf r) = \lambda^3
        n(\lambda \mathbf r)$. Then
        $E_x[n_\lambda] = \lambda E_x[n]$ and
        $E_c[n_\lambda] < \lambda E_c[n]$ for $\lambda > 1$.
  \item \textbf{Lieb--Oxford bound.}
        $E_x[n] \ge -C \int n(\mathbf r)^{4/3}\, d\mathbf r$ for a known
        constant $C$.
  \item \textbf{One-electron self-interaction freedom.}
        For any one-electron density,
        $\Exc[n] = -J[n]$, equivalently $E_x[n] = -J[n]$ and $E_c[n] = 0$.
        This is the constraint of central interest in this report.
  \item \textbf{Piecewise linearity in fractional particle number.}
        $E(N)$ is piecewise linear between integer $N$, with derivative
        discontinuities at integers.
  \item \textbf{Euclidean equivariance.}
        $\Exc$ is invariant under rotations, translations, and reflections
        of the density.
\end{itemize}

Functionals that are constructed to satisfy some of these constraints
analytically are called \emph{non-empirical}. PBE, SCAN, and r2SCAN are
examples. Functionals that are fitted to data with less attention to
constraints are called \emph{empirical}. B3LYP is a hybrid in this category.

\section{Self-interaction error and delocalization error}

\subsection{Self-interaction error (SIE)}

In exact DFT, an electron does not repel itself. The classical Coulomb
term $J[n]$ does include a spurious self-repulsion for each electron, and
the exact $\Exc$ cancels it exactly. Approximate functionals do not cancel
it exactly, and the residual is called the \emph{self-interaction error}.

For a system with exactly one electron, the true electron--electron energy
is zero (there is no second electron). Therefore the exact $\Exc$ must
cancel $J[n]$ exactly:
\begin{equation}
\Exc[n] = -J[n] \qquad (N = 1).
\label{eq:exact-constraint-n1}
\end{equation}
Any approximate functional that violates Eq.~\eqref{eq:exact-constraint-n1}
produces a non-zero \emph{one-electron SIE}. This is the quantity measured
in this work.

\subsection{Delocalization error}

Delocalization error is the tendency of approximate functionals to
over-stabilize charge distributions that spread electron density over
multiple centers, relative to the exact distribution. It is closely related
to SIE: incomplete cancellation of $J[n]$ makes concentrating charge on a
single center appear artificially unfavorable, so the functional prefers
spreading it. On a stretched bond, delocalization error manifests as an
artificial over-binding of the symmetric, delocalized solution.

\section{The L\"owdin symmetry dilemma}

Many approximate functionals have more than one self-consistent-field (SCF)
solution for a given nuclear geometry. When the true ground state is
symmetric under some operation (e.g., reflection through the bond midpoint
in H$_2^+$), but the approximate functional's lowest-energy SCF solution
breaks that symmetry, the functional suffers from the \emph{symmetry
dilemma} identified by L\"owdin in the 1960s.

For a one-electron bond such as H$_2^+$, the symmetry dilemma is
particularly stark: the exact ground state is symmetric at every bond
length, but approximate functionals can converge to a symmetry-broken
solution with the electron localized on one center. This is the phenomenon
this report is primarily about.

\section{Machine-learned exchange--correlation functionals}

A recent development is the use of machine learning to represent $\Exc$.
Two current-generation examples are:

\begin{itemize}
  \item \textbf{CIDER} (Bystrom \& Kozinsky). A family of functionals whose
        nonlocal descriptors are constructed to satisfy the
        uniform-coordinate-scaling constraint by design. The released
        model used in this work is \texttt{CIDER26XCCHEM} (2026
        generation).
  \item \textbf{Skala} (Microsoft Research, 2025). A neural-network
        functional built on \texttt{e3nn} layers, which are Euclidean
        ($E(3)$) equivariant by construction.
\end{itemize}

Both encode one exact constraint into the architecture, independent of what
is learned on top. Whether such architectural guarantees confer protection
against \emph{other} pathologies --- in particular the one-electron
symmetry dilemma --- is an open empirical question. This report answers it,
in one concrete case, for CIDER26XCCHEM.

% =====================================================================
\chapter{Theory}
\label{ch:theory}
% =====================================================================

This chapter states and derives the theoretical results that the
calculations rely on. Nothing here is new mathematics; the derivations are
included because the whole diagnostic framework depends on them and because
later chapters cite them repeatedly.

\section{Hartree--Fock is exact for one electron}

\subsection{The Hartree--Fock energy expression}

For an $N$-electron system described by a single Slater determinant built
from orthonormal spin-orbitals $\varphi_1, \dots, \varphi_N$, the
variational minimization of $\langle \hat H \rangle$ yields the standard
Hartree--Fock total energy
\begin{equation}
\EHF = \sum_{i=1}^N \langle \varphi_i | \hat h | \varphi_i \rangle
       + \tfrac12 \sum_{i,j=1}^N \big[ J_{ij} - K_{ij} \big],
\label{eq:hf-general}
\end{equation}
where
\begin{equation}
\hat h = -\tfrac12 \nabla^2 + v_{\rm ext}(\mathbf r)
\end{equation}
is the one-electron operator (kinetic energy plus nuclear attraction; it
contains no electron--electron term), and the two-electron integrals are
\begin{align}
J_{ij} &= \iint |\varphi_i(\mathbf r_1)|^2
                \frac{1}{|\mathbf r_1 - \mathbf r_2|}
                |\varphi_j(\mathbf r_2)|^2 \, d\mathbf r_1 \, d\mathbf r_2,
\label{eq:Jij}\\
K_{ij} &= \iint \varphi_i^*(\mathbf r_1) \varphi_j(\mathbf r_1)
                \frac{1}{|\mathbf r_1 - \mathbf r_2|}
                \varphi_j^*(\mathbf r_2) \varphi_i(\mathbf r_2)
                \, d\mathbf r_1 \, d\mathbf r_2.
\label{eq:Kij}
\end{align}
$K_{ij}$ is non-zero only when $\varphi_i$ and $\varphi_j$ have the same
spin; this detail is irrelevant below.

\subsection{Specializing to $N=1$}

With one occupied spin-orbital $\varphi_1$, the double sum in
Eq.~\eqref{eq:hf-general} collapses to its single $i=j=1$ term:
\begin{equation}
\EHF(N{=}1) = \langle \varphi_1 | \hat h | \varphi_1 \rangle
              + \tfrac12 \big[ J_{11} - K_{11} \big].
\label{eq:hf-n1}
\end{equation}

\begin{lemma}[$J_{11} = K_{11}$]
\label{lem:J11}
For any normalized spin-orbital $\varphi_1$, $J_{11} = K_{11}$ exactly.
\end{lemma}

\begin{proof}
Substituting $i=j=1$ in Eqs.~\eqref{eq:Jij} and~\eqref{eq:Kij},
\begin{align}
J_{11} &= \iint |\varphi_1(\mathbf r_1)|^2
                \frac{1}{|\mathbf r_1 - \mathbf r_2|}
                |\varphi_1(\mathbf r_2)|^2 \, d\mathbf r_1 \, d\mathbf r_2,\\
K_{11} &= \iint \varphi_1^*(\mathbf r_1) \varphi_1(\mathbf r_1)
                \frac{1}{|\mathbf r_1 - \mathbf r_2|}
                \varphi_1^*(\mathbf r_2) \varphi_1(\mathbf r_2)
                \, d\mathbf r_1 \, d\mathbf r_2.
\end{align}
But $\varphi_1^*(\mathbf r) \varphi_1(\mathbf r) = |\varphi_1(\mathbf r)|^2$
identically, and the same at $\mathbf r_2$. Substituting, the integrand of
$K_{11}$ becomes term-by-term identical to the integrand of $J_{11}$.
Hence $K_{11} = J_{11}$ for any normalized $\varphi_1$.
\end{proof}

\begin{theorem}[HF exactness for $N=1$]
\label{thm:hf-exact}
For any one-electron system in any one-particle basis,
\begin{equation}
\EHF(N{=}1) = \langle \varphi_1 | \hat h | \varphi_1 \rangle,
\end{equation}
with zero net electron--electron energy. This is the exact
(full-CI-in-that-basis) ground-state energy of $\hat h$ in that basis.
\end{theorem}

\begin{proof}
Substituting Lemma~\ref{lem:J11} into Eq.~\eqref{eq:hf-n1} gives
$\tfrac12 (J_{11} - K_{11}) = 0$, so
$\EHF(N{=}1) = \langle \varphi_1 | \hat h | \varphi_1 \rangle$.
A one-electron system has no $1/r_{12}$ term in its Hamiltonian, so the
true electron--electron energy is zero and the exact total energy is just
the expectation value of $\hat h$ in the ground state. Antisymmetrization
is vacuous for one particle, so a single Slater determinant places no
restriction on the wavefunction's form. Minimizing $\langle \varphi_1 |
\hat h | \varphi_1 \rangle$ over normalized $\varphi_1$ in a finite basis
is therefore identical to diagonalizing $\hat h$ in that basis and taking
the lowest eigenvalue.
\end{proof}

\subsection{The exact constraint on $\Exc$}

The classical Coulomb functional of a one-electron density
$n(\mathbf r) = |\varphi_1(\mathbf r)|^2$ is
\begin{equation}
J[n] = \tfrac12 \iint \frac{n(\mathbf r_1) n(\mathbf r_2)}
                           {|\mathbf r_1 - \mathbf r_2|}
       \, d\mathbf r_1 \, d\mathbf r_2 = J_{11}.
\end{equation}
Theorem~\ref{thm:hf-exact} says the true electron--electron energy is zero
for $N=1$, while $J[n] > 0$ for any non-point density. The only way both
statements can hold is
\begin{equation}
\Exc[n] = -J[n] \qquad (N=1).
\end{equation}
Correlation is defined as $E_c = \Eexact - \EHF$. Theorem~\ref{thm:hf-exact}
shows $\Eexact = \EHF$ for $N=1$, so $E_c[n] = 0$ and the entire
cancellation is carried by exchange:
\begin{equation}
\boxed{\; E_x[n] = -J[n], \quad E_c[n] = 0 \qquad (N=1). \;}
\label{eq:sic-constraint}
\end{equation}
This is the constraint first stated by Perdew and Zunger (1981).

\section{The same-basis reference}

\subsection{Why the naive reference contaminates SIE measurements}

A natural reference for a one-electron SIE measurement is the exact
hydrogenic energy $E = -Z^2/2$ (atomic units), which is the exact
infinite-basis ground-state energy of a hydrogen-like ion of nuclear charge
$Z$. However, any real calculation uses a finite one-particle basis
(Gaussian functions in this work). The finite basis cannot represent the
exact hydrogenic orbital perfectly, so it introduces a
\emph{basis-incompleteness} error separate from any functional error.

If the finite-basis SIE is measured against the infinite-basis value,
the result is
\begin{equation}
\SIE_{\rm naive} = E_{\rm DFT}^{\rm finite} - E_{\rm exact}^{\infty},
\end{equation}
which conflates two errors:
\begin{equation}
\SIE_{\rm naive}
= \underbrace{\big( E_{\rm DFT}^{\rm finite} - E_{\rm exact}^{\rm finite} \big)}_{\text{functional error}}
+ \underbrace{\big( E_{\rm exact}^{\rm finite} - E_{\rm exact}^{\infty} \big)}_{\text{basis error}}.
\end{equation}
The second term is not a property of the functional and contaminates the
diagnostic.

\subsection{The same-basis reference}

Theorem~\ref{thm:hf-exact} supplies a cleaner reference: run unrestricted
Hartree--Fock (UHF) in the \emph{same} finite basis as the DFT calculation.
By the theorem, the UHF energy is then \emph{exactly} the
full-CI-in-that-basis ground-state energy, with no approximation. The
same-basis SIE is
\begin{equation}
\SIE_{\rm same} = E_{\rm DFT}^{\rm finite} - \EHF^{\rm finite},
\end{equation}
and the basis-incompleteness terms cancel by construction. This is the
reference used throughout this report.

\begin{remark}
The same-basis UHF reference is exact for any one-electron system, at any
geometry, in any basis. In particular it applies to H$_2^+$ at \emph{every}
internuclear separation, not just at the equilibrium geometry.
\end{remark}

\section{Delocalization error on a stretched one-electron bond}

H$_2^+$ has one electron at every internuclear separation $R$. Its exact
ground state is symmetric under reflection through the bond midpoint, and
its exact ground-state density is symmetric, so the exact electron
population on each proton is $0.5$. This is Theorem~\ref{thm:h2-sym} below.

\subsection{Symmetry of the exact ground state}

The electronic Hamiltonian of H$_2^+$ (Born--Oppenheimer, protons fixed at
$\pm (R/2) \hat{\mathbf z}$) is
\begin{equation}
\hat H_{\rm el} = -\tfrac12 \nabla^2
                  - \frac{1}{|\mathbf r - \tfrac{R}{2} \hat{\mathbf z}|}
                  - \frac{1}{|\mathbf r + \tfrac{R}{2} \hat{\mathbf z}|}.
\end{equation}
Let $\hat \sigma$ denote reflection through the plane $z = 0$,
$\hat \sigma \psi(x,y,z) = \psi(x,y,-z)$.

\begin{lemma}
\label{lem:sym-commute}
$[\hat H_{\rm el}, \hat \sigma] = 0$.
\end{lemma}

\begin{proof}
$-\tfrac12 \nabla^2$ is invariant under any reflection. Applying
$\hat \sigma$ to the potential sends the two attraction terms into each
other (since $|(x,y,-z) - (R/2)\hat{\mathbf z}| = |\mathbf r +
(R/2)\hat{\mathbf z}|$), and their sum is unchanged.
\end{proof}

\begin{theorem}[Symmetric exact density]
\label{thm:h2-sym}
If the electronic ground state of H$_2^+$ is non-degenerate, its density
satisfies $n(x,y,z) = n(x,y,-z)$ for all $\mathbf r$ and all $R$.
Equivalently, the exact electron population on each proton is exactly
$0.5$ at every $R$.
\end{theorem}

\begin{proof}
$\hat \sigma$ is a reflection, so $\hat \sigma^2 = \hat 1$ and its
eigenvalues are $\pm 1$. By Lemma~\ref{lem:sym-commute}, $\hat H_{\rm el}$
and $\hat \sigma$ commute, so they share an eigenbasis. If $\psi_0$ is a
non-degenerate eigenstate of $\hat H_{\rm el}$ with eigenvalue $E_0$, then
$\hat \sigma \psi_0$ is also an eigenstate with the same eigenvalue, so by
non-degeneracy $\hat \sigma \psi_0 = \lambda \psi_0$ for some scalar
$\lambda$. Applying $\hat \sigma$ again gives $\lambda^2 = 1$, so
$\lambda = \pm 1$. Either way
$|\hat \sigma \psi_0|^2 = |\psi_0|^2$, so
$n(x,y,z) = n(x,y,-z)$. The lowest eigenvalue of a bound two-center
one-particle Hamiltonian of this form is the bonding $1\sigma_g$-like
state, non-degenerate at every $R$.
\end{proof}

\begin{corollary}[Exact Mulliken charges]
For the exact ground state of H$_2^+$, the Mulliken charge on each proton
is exactly $q(\mathrm H_a) = q(\mathrm H_b) = 0.5$ at every $R$.
\end{corollary}

\begin{proof}
Mulliken population assigns $p_A = \sum_{\mu \in A} \sum_\nu D_{\mu \nu}
S_{\nu \mu}$ to atom $A$. By Theorem~\ref{thm:h2-sym}, the exact density is
symmetric under $\hat \sigma$, and the two centers are exchanged by
$\hat \sigma$ with identical basis sets, so $p_a = p_b$. Since
$p_a + p_b = 1$ (one electron total) and $Z_a = Z_b = 1$, this gives
$q_a = q_b = 0.5$.
\end{proof}

Any numerical calculation returning $q_a \ne q_b$ on H$_2^+$ is, by
Theorem~\ref{thm:h2-sym}, not describing the true ground-state density.
It has converged to a different critical point of the approximate energy
functional. This is L\"owdin's symmetry dilemma.

\subsection{Definition of the delocalization error}

Define the signed delocalization error at bond length $R$ as
\begin{equation}
\Delta_{\rm deloc}(R) = \ECD(R) - \EHF(N{=}1; R),
\label{eq:deloc-def}
\end{equation}
where both terms use the same finite basis. Since $\EHF$ is exact by
Theorem~\ref{thm:hf-exact}, $\Delta_{\rm deloc}$ is the pure functional
error at that geometry.

Sign convention: $\Delta_{\rm deloc} < 0$ means the functional's energy is
below the exact reference, i.e., the functional \emph{over-stabilizes}
the system. On a stretched one-electron bond, this is the signature of
delocalization error. $\Delta_{\rm deloc} > 0$ means the functional
\emph{under}-stabilizes, which is a different pathology.

% =====================================================================
\chapter{Methods}
\label{ch:methods}
% =====================================================================

\section{Software and hardware}

All calculations use PySCF. The machine-learned functional
\texttt{CIDER26XCCHEM} is loaded via the open-source \texttt{ciderpress}
package, whose pretrained weights ship inside the pip package. No external
network access is required after installation.

\todo{Record the exact PySCF version, ciderpress version, and Python version
used for every calculation. This is required for reproducibility.}

\section{Basis sets}

The primary basis is \texttt{def2-TZVP}, a triple-$\zeta$ valence basis
with polarization functions. It is large enough that basis-incompleteness
effects are small but not negligible at the millihartree level relevant to
this work (see Section~\ref{sec:basis-incompleteness}).

\todo{Add calculations with \texttt{def2-SVP}, \texttt{def2-QZVP}, and at
least one \texttt{aug-cc-pVXZ} basis to demonstrate that conclusions are
not basis-set specific. This is currently missing.}

\section{SCF settings}

For DFT calculations, unrestricted Kohn--Sham (\texttt{dft.UKS}) is used.
For the exact reference, unrestricted Hartree--Fock (\texttt{scf.UHF}) is
used. Tight convergence settings:
\begin{itemize}
  \item energy tolerance: $10^{-10}$ Hartree,
  \item maximum cycles: 200,
  \item default initial guess unless otherwise noted.
\end{itemize}

\todo{Every SCF calculation must record: convergence flag, number of cycles,
final energy, and (for symmetry-relevant systems) Mulliken charges. In
addition, an SCF stability analysis (\texttt{mf.stability()}) must be
performed at every geometry for every functional whose symmetry-breaking
behavior is being studied. This was missing from D2 and D3 as originally
drafted.}

\section{Systems studied}

\subsection{Isoelectronic series}

H ($Z = 1$), He$^+$ ($Z = 2$), Li$^{2+}$ ($Z = 3$), Be$^{3+}$ ($Z = 4$).
These are single atoms, so no bond-length scan is needed.

\subsection{H$_2^+$ dissociation}

14 bond lengths in the original D3 scan:
$R \in \{0.5, 0.75, 1.06, 1.5, 2.0, 2.5, 3.0, 3.5, 4.0, 4.5, 5.0, 7.0,
10.0, 15.0\}$~\AA. Two additional points (5.5, 6.0) are needed to resolve
the bifurcation location.

\todo{Add $R = 5.5$ and $R = 6.0$~\AA.}

\subsection{Regression set}

For the blast-radius check (Section~\ref{sec:patch}), four closed-shell
molecules: H$_2$O, CH$_4$, N$_2$, Ne.

\section{The master dataset}

All results are collected into a single CSV file with the following
columns:

\begin{center}
\begin{tabular}{ll}
\toprule
Column & Meaning \\
\midrule
\texttt{functional} & PBE, r2SCAN, B3LYP, CIDER23X, CIDER26XCCHEM, or UHF \\
\texttt{system} & ion name or molecule \\
\texttt{R} & bond length in \AA\ (blank for atoms) \\
\texttt{E} & total energy in Hartree \\
\texttt{converged} & boolean \\
\texttt{cycles} & number of SCF cycles \\
\texttt{stable} & boolean (from \texttt{mf.stability()}) \\
\texttt{q\_Ha} & Mulliken charge on first atom \\
\texttt{q\_Hb} & Mulliken charge on second atom (blank for atoms) \\
\bottomrule
\end{tabular}
\end{center}

Every number in every table and figure in this report must be traceable to
a row of this CSV. No hand-typed numbers anywhere.

\todo{Build the master CSV from scratch. The numbers currently in D1--D5
are inconsistent across drafts (see Chapter~\ref{ch:issues}).}

% =====================================================================
\chapter{Results}
\label{ch:results}
% =====================================================================

\section{Basis-incompleteness of the same-basis reference}
\label{sec:basis-incompleteness}

Table~\ref{tab:basis} compares the same-basis UHF energy of the
isoelectronic series against the exact infinite-basis hydrogenic value
$-Z^2/2$. The difference is the basis-incompleteness term, which is what
contaminates the naive SIE measurement.

\begin{table}[htbp]
\centering
\caption{Same-basis UHF energies vs.\ infinite-basis hydrogenic values.
Basis: \texttt{def2-TZVP}. Incompleteness = $E_{\rm UHF} - (-Z^2/2)$.}
\label{tab:basis}
\begin{tabular}{lrrrr}
\toprule
Ion & $Z$ & $-Z^2/2$ (Ha) & $E_{\rm UHF}$ same-basis (Ha) & Incompleteness (mHa) \\
\midrule
H        & 1 & $-0.500000$ & $-0.499810$ & $0.190$ \\
He$^+$   & 2 & $-2.000000$ & $-1.998140$ & $1.860$ \\
Li$^{2+}$& 3 & $-4.500000$ & $-4.498572$ & $1.428$ \\
Be$^{3+}$& 4 & $-8.000000$ & $-7.996125$ & $3.875$ \\
\bottomrule
\end{tabular}
\end{table}

Two observations:
\begin{enumerate}
  \item The incompleteness is small in absolute terms ($0.2$--$4$~mHa) but
        comparable to the SIE residuals being measured (also a few mHa).
  \item The incompleteness is \emph{not} monotonic in $Z$: He$^+$ has
        larger incompleteness than Li$^{2+}$. This contradicts the narrative
        in D1, which claimed monotonic growth. The correct statement is
        that the incompleteness is largest for Be$^{3+}$ among these four
        ions, but it is not a monotonic function of $Z$ for this basis.
\end{enumerate}

\todo{This table must be recomputed and re-verified. The non-monotonicity
must be reported honestly, not hidden.}

\section{Reproduction on known functionals}

Table~\ref{tab:repro} reports the same-basis SIE on the H atom and the
same-basis delocalization error on stretched H$_2^+$ at $R = 10$~\AA, for
three functionals whose behavior is well-established in the literature.
The purpose of this section is to demonstrate that the diagnostic
reproduces known rankings before it is applied to a new functional.

\begin{table}[htbp]
\centering
\caption{Same-basis SIE (H atom) and $\Delta_{\rm deloc}$ (H$_2^+$,
$R = 10$~\AA) for three known functionals. Basis: \texttt{def2-TZVP}.
Values in mHa.}
\label{tab:repro}
\begin{tabular}{lrrr}
\toprule
Functional & H-atom SIE & H$_2^+$ $R{=}10$~\AA\ $\Delta$ & Comment \\
\midrule
PBE    & $+0.384$ & $+0.798$  & GGA \\
r2SCAN & $+0.073$ & $+0.114$  & meta-GGA, constraint-aware \\
B3LYP  & $-2.154$ & $-77.893$ & 20\% hybrid \\
\bottomrule
\end{tabular}
\end{table}

\todo{The B3LYP number $-77.893$~mHa is implausibly large. At
$R = 10$~\AA, the exact separated-fragments limit is
$E \approx -0.5$~Ha. A B3LYP energy of $-0.577893$~Ha at this geometry
implies the SCF has converged to a solution roughly $78$~mHa below the
exact separated-fragments energy, which is either (a) a symmetry-broken
solution with the electron localized on one proton (in which case the
energy should be near $-0.5$~Ha, not $-0.578$~Ha, and the number is
wrong), (b) a non-converged SCF that hit the cycle limit, or (c) a bug
in the input. Before this number is used anywhere, the calculation must
be rerun with SCF stability analysis and Mulliken charge diagnostics
reported.}

\todo{The PBE and r2SCAN signs at $R = 10$~\AA\ are \emph{positive},
meaning these functionals under-bind the stretched bond rather than
over-bind it. The narrative in D2 called this ``delocalization error,''
which is incorrect: delocalization error is over-binding
($\Delta < 0$). Either the sign convention is inverted somewhere, or PBE
and r2SCAN are showing a different pathology (under-binding) that requires
separate framing. This must be resolved before any claim is made about
these two functionals.}

\section{CIDER26XCCHEM on the isoelectronic series}

\begin{table}[htbp]
\centering
\caption{CIDER26XCCHEM SIE on the isoelectronic series, same-basis
reference. Basis: \texttt{def2-TZVP}. Values in mHa.}
\label{tab:cider-iso}
\begin{tabular}{lrr}
\toprule
Ion & Naive SIE (vs.\ $-Z^2/2$) & Same-basis SIE (vs.\ UHF) \\
\midrule
H        & $0.574$ & $0.383$ \\
He$^+$   & $5.047$ & $3.187$ \\
Li$^{2+}$& $5.879$ & $4.451$ \\
Be$^{3+}$& $9.784$ & $5.909$ \\
\bottomrule
\end{tabular}
\end{table}

The naive reference overestimates the same-basis SIE by a ratio of
roughly $1.32$--$1.66$ across this series. This is the quantitative
content of D1's recommendation: same-basis reference is not a cosmetic
correction but changes the reported SIE by 30--66\%.

\todo{These numbers must be recomputed as part of the master dataset.
Cross-paper inconsistencies exist (see Chapter~\ref{ch:issues}).}

\section{CIDER26XCCHEM on the H$_2^+$ dissociation curve}

\subsection{Data}

Table~\ref{tab:h2plus} reproduces the D3 H$_2^+$ scan. Columns:
$R$; CIDER26XCCHEM energy; SCF convergence flag; Mulliken charges
$q(\mathrm H_a)$, $q(\mathrm H_b)$; same-basis UHF reference energy;
signed error $\Delta = \ECD - \EHF$; branch assignment.

\begin{longtable}{rrrrrrl}
\caption{CIDER26XCCHEM on the H$_2^+$ dissociation curve. Basis:
\texttt{def2-TZVP}; convergence tolerance $10^{-10}$~Ha, up to 200 cycles.
\label{tab:h2plus}}\\
\toprule
$R$ (\AA) & $\ECD$ (Ha) & conv.\ & $q_a$ & $q_b$ & $\EHF$ (Ha) & $\Delta$ (mHa) \\
\midrule
\endfirsthead
\toprule
$R$ (\AA) & $\ECD$ (Ha) & conv.\ & $q_a$ & $q_b$ & $\EHF$ (Ha) & $\Delta$ (mHa) \\
\midrule
\endhead
0.50  & $-0.417745$ & yes & 0.500 & 0.500 & $-0.417839$ & $+0.094$ \\
0.75  & $-0.571268$ & yes & 0.500 & 0.500 & $-0.571108$ & $-0.160$ \\
1.06  & $-0.601529$ & yes & 0.500 & 0.500 & $-0.600965$ & $-0.564$ \\
1.50  & $-0.579842$ & yes & 0.500 & 0.500 & $-0.578759$ & $-1.083$ \\
2.00  & $-0.548119$ & yes & 0.500 & 0.500 & $-0.545618$ & $-2.502$ \\
2.50  & $-0.533070$ & yes & 0.500 & 0.500 & $-0.524201$ & $-8.869$ \\
3.00  & $-0.528867$ & yes & 0.500 & 0.500 & $-0.512060$ & $-16.808$ \\
3.50  & $-0.529214$ & yes & 0.500 & 0.500 & $-0.505617$ & $-23.597$ \\
4.00  & $-0.531640$ & yes & 0.500 & 0.500 & $-0.502439$ & $-29.201$ \\
4.50  & $-0.534476$ & yes & 0.500 & 0.500 & $-0.500961$ & $-33.515$ \\
5.00  & $-0.499066$ & no  & 0.000 & 1.000 & $-0.500305$ & $+1.238$ \\
7.00  & $-0.498653$ & no  & 1.000 & 0.000 & $-0.499828$ & $+1.174$ \\
10.00 & $-0.499350$ & no  & 1.000 & 0.000 & $-0.499812$ & $+0.462$ \\
15.00 & $-0.499351$ & no  & 0.000 & 1.000 & $-0.499810$ & $+0.459$ \\
\bottomrule
\end{longtable}

\subsection{Reading the curve}

The data fall into three regimes:

\begin{enumerate}
  \item \textbf{Near equilibrium} ($R = 0.5$--$1.5$~\AA): SCF converges to
        the symmetric solution ($q_a = q_b = 0.5$), error is small
        ($|\Delta| < 1.1$~mHa).
  \item \textbf{Stretching regime} ($R = 2.0$--$4.5$~\AA): SCF still
        converges to the symmetric solution, but the error grows
        monotonically from $-2.5$~mHa to $-33.5$~mHa. The negative sign
        means the functional over-stabilizes the symmetric solution, which
        is textbook delocalization error.
  \item \textbf{Beyond the bifurcation} ($R \ge 5.0$~\AA): SCF does
        \emph{not} converge within 200 cycles at $10^{-10}$ tolerance, and
        the Mulliken charges are either $(0,1)$ or $(1,0)$, i.e., the
        electron is localized on one proton. The apparent error is small
        ($|\Delta| < 1.3$~mHa), but this is not a recovery; the functional
        has switched to describing a different physical state (separated
        fragments) rather than the true symmetric ground state.
\end{enumerate}

\subsection{The bifurcation}

Between $R = 4.5$~\AA\ and $R = 5.0$~\AA, the character of the SCF
solution changes qualitatively. This is a genuine bifurcation of the SCF
solution landscape. \todo{It must be confirmed by:}

\begin{itemize}
  \item \todo{SCF stability analysis (\texttt{mf.stability()}) at every
        geometry, to confirm both the symmetric and localized solutions
        are genuine SCF stationary points.}
  \item \todo{A second SCF calculation at each $R \ge 5.0$~\AA\ starting
        from a symmetric initial guess, to confirm the symmetric solution
        still exists (even if it is not the lowest one).}
  \item \todo{Mulliken charges reported for both solutions.}
\end{itemize}

\subsection{Why the post-bifurcation error is small}

At $R \ge 5.0$~\AA, the localized solution describes a neutral H atom plus
a bare proton at large separation. Its energy is close to the isolated
H-atom energy, which has a small SIE ($\sim 0.4$~mHa for CIDER26XCCHEM,
from Table~\ref{tab:cider-iso}). The small apparent error beyond the
bifurcation is therefore not a sign of the functional doing better; it is a
sign of the functional having stopped attempting to describe the true
delocalized state.

\section{Refutation of the one-electron patch hypothesis}
\label{sec:patch}

\subsection{The hypothesis}

Hypothesis H1 (from D4): if CIDER26XCCHEM's learned enhancement map is
constrained to satisfy the exact one-electron limit, its SIE should shrink
toward r2SCAN's without degrading multi-electron accuracy.

\subsection{The proposed test and why it is vacuous}

The most direct test is a global, per-system energy substitution gated on
electron count:
\begin{equation}
E_{\rm patched}[\text{system}] =
\begin{cases}
\EHF(N{=}1;\text{same basis}) & \text{if } N_{\rm electron} = 1, \\
E_\phi[\text{system}] & \text{otherwise}.
\end{cases}
\label{eq:patch}
\end{equation}

By Theorem~\ref{thm:hf-exact}, $\EHF(N{=}1)$ is exact in any basis, so
the patched functional returns the exact answer on every one-electron
system, regardless of which underlying functional $\phi$ is wrapped. The
value $E_\phi$ is never consulted on the patched branch.

This means Eq.~\eqref{eq:patch} cannot discriminate between functionals:
PBE, B3LYP, and CIDER26XCCHEM all collapse to identical, exact results
when patched. The patch is tautologically correct, and running it provides
no evidence about the underlying functional's treatment of the one-electron
limit.

\begin{table}[htbp]
\centering
\caption{Raw vs.\ patched one-electron errors (mHa). The patched column
is identically zero for every functional, by construction.}
\label{tab:patch}
\begin{tabular}{lrrrrrr}
\toprule
Ion & PBE raw & PBE patched & B3LYP raw & B3LYP patched & CIDER raw & CIDER patched \\
\midrule
H        & $+0.194$ & $0$ & $-2.344$ & $0$ & $+0.384$ & $0$ \\
He$^+$   & $+6.050$ & $0$ & $+1.118$ & $0$ & $+3.187$ & $0$ \\
Li$^{2+}$& $+12.462$& $0$ & $+5.493$ & $0$ & $+4.451$ & $0$ \\
Be$^{3+}$& $+17.094$& $0$ & $+8.923$ & $0$ & $+5.909$ & $0$ \\
\bottomrule
\end{tabular}
\end{table}

\todo{The PBE and B3LYP raw numbers in this table must be reconciled with
Table~\ref{tab:repro} (they disagree). Both must come from the master CSV.}

\subsection{What a meaningful test would require}

A genuine test of H1 must modify the functional's \emph{internal
representation} --- e.g., its descriptor pipeline or its enhancement-factor
spline --- so that the functional's own one-electron prediction changes,
rather than being replaced after the fact. This is stated as a direction
for future work, not as a result of the present one.

\section{The constraint-residual poset is not totally ordered}

\subsection{The two-axis comparison}

Order a set of functionals by their residual vectors across multiple
constraint axes. Two functionals are \emph{incomparable} if each is smaller
on some axis. This is a partial order, not a total order.

Table~\ref{tab:poset} reports residuals on two axes: same-basis SIE on
the H atom, and same-basis delocalization error on H$_2^+$ at
$R = 10$~\AA, for five functionals.

\begin{table}[htbp]
\centering
\caption{Two-axis constraint residuals (mHa, absolute value).}
\label{tab:poset}
\begin{tabular}{lrr}
\toprule
Functional & $|\SIE|$ H atom & $|\Delta_{\rm deloc}|$ H$_2^+$ $R{=}10$~\AA \\
\midrule
PBE              & $0.384$ & $0.798$ \\
r2SCAN           & $0.073$ & $0.110$ \\
B3LYP            & $2.154$ & $77.893$ \\
CIDER23X         & $2.554$ & $2.412$ \\
CIDER26XCCHEM    & $0.574$ & $0.641$ \\
\bottomrule
\end{tabular}
\end{table}

Two pairs are incomparable under this order:
\begin{itemize}
  \item PBE vs.\ CIDER26XCCHEM: PBE has smaller SIE, CIDER has smaller
        delocalization error.
  \item B3LYP vs.\ CIDER23X: B3LYP has smaller SIE, CIDER23X has far
        smaller delocalization error.
\end{itemize}

\todo{These numbers inherit the B3LYP anomaly from
Section~\ref{sec:basis-incompleteness}. The B3LYP delocalization number
($77.893$~mHa) must be verified before the incomparability claim is used.
If B3LYP's number is wrong, the second incomparable pair may not survive.}

\subsection{What this shows}

A single scalar accuracy metric cannot summarize functional quality across
these two axes. A functional can be better on one and worse on another,
and no weighted combination of the two residuals can resolve this without
an arbitrary choice of weights.

% =====================================================================
\chapter{Claims inventory}
\label{ch:claims}
% =====================================================================

This chapter lists every claim the work makes, with an explicit
confidence tag. Claims are the load-bearing statements; everything else in
this report is scaffolding.

\section{Verified claims}

\begin{enumerate}
  \item \verified{Hartree--Fock theory is exact for any one-electron
        system, in any finite one-particle basis, at any geometry.}
        Proof: Theorem~\ref{thm:hf-exact}. This is textbook.
  \item \verified{The exact one-electron constraint on $\Exc$ is
        $E_x[n] = -J[n]$, $E_c[n] = 0$.} Proof:
        Eq.~\eqref{eq:sic-constraint}. Known since Perdew--Zunger 1981.
  \item \verified{Same-basis UHF is therefore the exact reference for
        one-electron SIE diagnostics, and comparing finite-basis DFT to
        the infinite-basis hydrogenic value contaminates the measurement.}
        This is a methodological recommendation; it is correct by
        construction.
  \item \verified{The exact ground state of H$_2^+$ is symmetric under
        reflection through the bond midpoint, at every $R$, and the exact
        Mulliken charges on the two protons are exactly $0.5$ each.}
        Proof: Theorem~\ref{thm:h2-sym}. This is textbook.
  \item \verified{A global, per-system energy substitution gated on
        electron count is tautologically exact on every one-electron
        system, for any underlying functional, and cannot discriminate
        between functionals.} Proof: Theorem~\ref{thm:hf-exact} applied to
        Eq.~\eqref{eq:patch}.
  \item \verified{On the two-axis constraint-residual space, PBE and
        CIDER26XCCHEM are incomparable.} Depends on
        Table~\ref{tab:poset} numbers; the incomparability argument is
        correct if the numbers are.
\end{enumerate}

\section{Tentative claims}

\begin{enumerate}
  \item \tentative{The naive infinite-basis reference inflates the
        apparent SIE of CIDER26XCCHEM on the isoelectronic series by
        30--66\%.} Based on Table~\ref{tab:cider-iso}; needs recomputation
        from the master CSV.
  \item \tentative{CIDER26XCCHEM develops a delocalization error on the
        symmetric H$_2^+$ branch that grows monotonically from
        $-2.5$~mHa at $R = 2.0$~\AA\ to $-33.5$~mHa at $R = 4.5$~\AA.}
        Based on Table~\ref{tab:h2plus}; needs SCF stability confirmation.
  \item \tentative{The H$_2^+$ SCF solution bifurcates between
        $R = 4.5$~\AA\ and $R = 5.0$~\AA\ to a symmetry-broken solution
        with the electron localized on one proton.} Based on Mulliken
        charges in Table~\ref{tab:h2plus}; needs stability analysis and
        a localized-initial-guess cross-check.
  \item \tentative{The apparent small error beyond the bifurcation is not
        a recovery but a signature of the functional describing separated
        fragments.} Physical argument is sound; needs quantitative
        confirmation (e.g., compare post-bifurcation energy to isolated
        H-atom SIE).
  \item \tentative{The constraint-residual order across the five
        functionals tested is not total.} Based on Table~\ref{tab:poset};
        depends on B3LYP number being correct.
\end{enumerate}

\section{Withdrawn or doubtful claims}

\begin{enumerate}
  \item \withdrawn{B3LYP shows a $77.9$~mHa delocalization error on
        H$_2^+$ at $R = 10$~\AA.} The number is implausibly large and
        likely reflects a non-converged or symmetry-broken SCF solution,
        not delocalization error. Must be diagnosed.
  \item \withdrawn{PBE and r2SCAN show delocalization error on stretched
        H$_2^+$ at $R = 10$~\AA.} Their signed errors are \emph{positive}
        (under-binding), not negative, so the label ``delocalization
        error'' is wrong for these two. The correct description may be
        under-binding of the symmetric solution; this needs separate
        framing.
  \item \withdrawn{The basis-incompleteness of the same-basis reference
        grows monotonically with $Z$.} Table~\ref{tab:basis} shows it is
        not monotonic (He$^+$ > Li$^{2+}$).
  \item \withdrawn{Skala and CIDER converge on the same abstract design
        principle.} D5 claimed this, but Skala was not actually run. The
        claim rests on the published architecture of Skala, not on
        independent verification.
  \item \withdrawn{The categorical framing of constraints as a poset
        provides new insight.} It provides bookkeeping, not new
        mathematics. The useful content is the two-axis incomparability
        argument, which does not require category theory.
\end{enumerate}

\section{Claims not made}

Explicitly not claimed in this work:

\begin{itemize}
  \item That CIDER26XCCHEM is worse than any other functional overall.
  \item That the bifurcation observed on H$_2^+$ occurs on any other
        system (multi-electron molecules, solids, surfaces).
  \item That the same-basis reference is the only correct reference
        (it is correct, but other references may also be correct for
        specific purposes).
  \item That a patch based on one-electron exactness is useless in
        general (it is useless for testing the hypothesis H1, but may
        have other uses).
  \item That Skala's architecture does or does not confer protection
        against delocalization error (untested).
\end{itemize}

% =====================================================================
\chapter{Known problems and open questions}
\label{ch:issues}
% =====================================================================

This chapter is the most important for deciding what to do next. Every
known problem with the work is listed here, with an assessment of severity
and a suggested resolution.

\section{Technical problems}

\subsection{B3LYP anomaly on stretched H$_2^+$}
\label{sec:b3lyp-problem}

\textbf{Severity: high.}

The D2 and D5 tables report B3LYP at $R = 10$~\AA\ giving
$E = -0.577893$~Ha, which is $78$~mHa below the exact separated-fragments
limit. This is either a non-converged SCF, a symmetry-broken solution, or
a bug. Until it is diagnosed, all claims depending on this number are
unsafe.

\textbf{Resolution:} rerun B3LYP on the H$_2^+$ curve with stability
analysis and Mulliken charges at every geometry. Report all solutions
found.

\subsection{Sign convention for PBE and r2SCAN delocalization error}

\textbf{Severity: high.}

D2's narrative calls the positive signed errors of PBE and r2SCAN
``delocalization error,'' but positive $\Delta$ means under-binding, not
over-binding. The terminology is wrong.

\textbf{Resolution:} either invert the sign convention everywhere (and
state it explicitly), or relabel the PBE/r2SCAN observations as
under-binding, distinct from delocalization.

\subsection{CIDER26XCCHEM SCF stability on H$_2^+$}

\textbf{Severity: high.}

The post-bifurcation ``localized branch'' energies were computed with
\texttt{converged=False}. A non-converged SCF energy is not guaranteed to
be a stationary point of the functional. The bifurcation might be an
artifact of a limit cycle, not a genuine second solution.

\textbf{Resolution:} run \texttt{mf.stability()} at every geometry;
confirm the localized solution is a genuine minimum. Run a second SCF
from a symmetric initial guess at $R \ge 5$~\AA\ to confirm the symmetric
solution still exists (as a saddle or higher minimum).

\subsection{Cross-paper numerical inconsistencies}

\textbf{Severity: high.}

The CIDER26XCCHEM H-atom SIE appears as $0.383$ in D1 and $0.574$ in D5.
The CIDER26XCCHEM H$_2^+$ $R = 10$~\AA\ $\Delta$ appears as $0.462$ in D3
and $0.641$ in D5. These are different numbers for what should be the same
quantity.

\textbf{Resolution:} rebuild everything from the master CSV. No
hand-typed numbers in any table or figure.

\subsection{Basis-set incompleteness not converged}

\textbf{Severity: medium.}

All results use \texttt{def2-TZVP}. Whether conclusions are basis-set
independent is untested.

\textbf{Resolution:} add \texttt{def2-SVP}, \texttt{def2-QZVP}, and
\texttt{aug-cc-pVTZ} runs on the key systems.

\subsection{Functional comparison incomplete}

\textbf{Severity: medium.}

Only PBE, r2SCAN, B3LYP, CIDER23X, and CIDER26XCCHEM are tested. SCAN is
missing.

\textbf{Resolution:} add SCAN. Optionally add a modern range-separated
hybrid (e.g., $\omega$B97X-V) as a positive control.

\subsection{The Skala claim is unverified}

\textbf{Severity: medium.}

D5 claims Skala and CIDER converge on the same design principle, but
Skala was not run.

\textbf{Resolution:} either run Skala (or a public reimplementation), or
downgrade the claim to a Discussion remark with explicit acknowledgment
that Skala was not tested.

\section{Presentation problems}

\subsection{AI coauthorship}

\textbf{Severity: critical for publication.}

All five drafts list ``Claude (Anthropic)'' as a coauthor. AI systems
cannot be authors under essentially all journal and arXiv policies.

\textbf{Resolution:} remove the AI coauthor from all author lists and the
bibliography. Add an acknowledgements line: ``The author used an AI
assistant for language editing and code drafting; all technical content
and calculations were verified by the author.''

\subsection{AI-slop language patterns}

\textbf{Severity: high for publication.}

Recurring patterns across the drafts:
\begin{itemize}
  \item ``to our knowledge for the first time''
  \item ``we believe,'' ``we think''
  \item ``precisely,'' ``exactly,'' ``genuinely,'' ``essentially'' used
        as intensifiers
  \item self-referential meta-commentary (``this is the crux of the
        paper,'' ``we record that correction here for the record'')
  \item defensive sentences (``this is not a numerical artifact,''
        ``we flag this honestly'')
  \item over-formalization of trivial results as Theorems/Corollaries
  \item grand verbs (``culminating in,'' ``systematized as,'' ``reproduced
        here in running code'')
\end{itemize}

\textbf{Resolution:} rewrite prose from scratch in the author's own voice.
Do not edit the existing drafts; start with a blank page and the master
CSV and figures.

\subsection{Circular citations}

\textbf{Severity: high for publication.}

Every draft cites the others as ``companion manuscript (2026),'' creating
mutual dependencies that no referee can evaluate.

\textbf{Resolution:} consolidate into one paper. Do not submit as five
mutually-citing manuscripts.

\subsection{Over-formalization}

\textbf{Severity: medium.}

The $J_{11} = K_{11}$ identity does not need a Lemma and Proof. The patch
refutation does not need a Theorem and Corollary. The categorical framework
of D5 does not need a category-theoretic construction.

\textbf{Resolution:} state trivial results as one-sentence observations.
Reserve formal environments for genuinely non-trivial results (of which
this work has essentially none; the real content is numerical).

\section{Open scientific questions}

\begin{enumerate}
  \item Does the H$_2^+$ bifurcation occur with other machine-learned
        functionals? With other one-electron bonds (HeH$^{2+}$,
        H$_2^{2+}$)?
  \item Does the bifurcation occur on any multi-electron system, or is
        it specific to the one-electron case?
  \item What is the relationship between the bifurcation and the
        functional's violation of piecewise linearity in fractional
        particle number? D3's Section 5 speculated on this but did not
        demonstrate it.
  \item Does CIDER26XCCHEM's scaling-covariance guarantee (from its
        descriptor design) provide any protection against
        delocalization error? The H$_2^+$ result suggests it does not,
        but this is one system.
  \item What is the right architecture for a machine-learned functional
        that avoids the bifurcation? This is the natural follow-up
        question, but it is a separate research project.
\end{enumerate}

% =====================================================================
\chapter{Future work}
\label{ch:future}
% =====================================================================

\section{Immediate actions (next two weeks)}

In order:

\begin{enumerate}
  \item \textbf{Diagnose B3LYP on stretched H$_2^+$.} Rerun with
        \texttt{mf.stability()} and Mulliken charges at every geometry.
        Determine whether the $-78$~mHa is a non-converged SCF, a
        symmetry-broken solution, or a bug. This is the highest-priority
        calculation because it determines whether D2 and D5 retain any
        content.

  \item \textbf{Confirm CIDER26XCCHEM bifurcation.} Run
        \texttt{mf.stability()} at every $R$ in the H$_2^+$ scan. Add
        $R = 5.5$ and $R = 6.0$~\AA. Run a second SCF at $R \ge 5$~\AA\
        from a symmetric initial guess. Confirm the localized solution is
        a genuine minimum.

  \item \textbf{Build the master CSV.} One file, one convention, all
        functionals, all systems. Every subsequent table and figure is
        generated from this file.

  \item \textbf{Recompute all numbers.} Regenerate every table and figure
        from the master CSV. Do not carry any hand-typed number forward.
\end{enumerate}

\section{Short-term actions (next two months)}

\begin{enumerate}
  \item Add \texttt{def2-SVP}, \texttt{def2-QZVP} basis sets to key
        systems to demonstrate basis independence.
  \item Add SCAN and one range-separated hybrid to the comparison.
  \item Run Skala if possible; otherwise remove the Skala claim.
  \item Write the paper (see Section~\ref{sec:paper-plan}).
\end{enumerate}

\section{Long-term questions (next year)}

\begin{enumerate}
  \item Extend the diagnostic to multi-electron systems. The
        same-basis UHF reference is exact only for $N = 1$; for
        $N > 1$ a different exact reference is needed (e.g., full CI in
        the same basis). This is computationally expensive but
        conceptually straightforward.
  \item Extend to excited states and to non-equilibrium geometries.
  \item Extend to periodic systems (solid-state DFT). The H$_2^+$
        pathology likely has a solid-state analogue in stretched
        one-electron bonds.
  \item Investigate the relationship between the bifurcation and the
        fractional-charge piecewise-linearity constraint, which D3
        speculated about but did not demonstrate.
\end{enumerate}

\section{Recommended paper structure}
\label{sec:paper-plan}

If the CIDER26XCCHEM bifurcation survives stability analysis, write
\textbf{one paper}, not five. Structure:

\begin{enumerate}
  \item \textbf{Introduction} (1 page). DFT, exact constraints,
        one-electron SIE, delocalization error, machine-learned
        functionals. State that the paper tests whether an
        architecturally constraint-aware ML functional is protected
        against the one-electron symmetry dilemma.

  \item \textbf{Theory and methods} (1.5 pages). Same-basis HF exactness
        (Theorem~\ref{thm:hf-exact}, stated but not over-formalized),
        same-basis reference, definition of $\SIE$ and $\Delta_{\rm
        deloc}$, computational setup.

  \item \textbf{Validation} (1 page). Reproductions on PBE, r2SCAN,
        B3LYP on the isoelectronic series and on H$_2^+$ at
        $R = 10$~\AA. Show that the diagnostic recovers known rankings.

  \item \textbf{Main result: H$_2^+$ bifurcation} (2 pages). The
        CIDER26XCCHEM H$_2^+$ curve, symmetric branch, delocalization
        error growth, bifurcation, localized branch, stability analysis.

  \item \textbf{Discussion} (1 page). Why the apparent post-bifurcation
        error recovery is not a recovery. The two-axis constraint-residual
        poset and the non-totality of functional ordering. Brief remark
        on the patch refutation.

  \item \textbf{Conclusion} (0.5 page). Recommendation: same-basis UHF
        as standard reference for one-electron SIE diagnostics; ML
        functional designers should test for symmetry breaking on
        stretched one-electron bonds.

  \item \textbf{Supporting information}: code, master CSV, convergence
        tables, additional basis sets.
\end{enumerate}

Target journal: \emph{J. Chem. Theory Comput.} or \emph{J. Chem. Phys.}
Length: 6--8 pages main text plus SI.

If the bifurcation does \emph{not} survive stability analysis, the paper
shrinks to a methods note on the same-basis reference, probably not
publishable in a top journal but suitable for a pedagogical venue or as
a section in someone else's paper.

% =====================================================================
\appendix
% =====================================================================

\chapter{Code snippets}

\section{PySCF setup for one-electron calculations}

\begin{verbatim}
from pyscf import gto, dft, scf

def run_dft(system, xc, R=None, basis="def2-tzvp", tight=True):
    if R is None:
        atom = system
        charge, spin = 0, 1   # adjust per ion
    else:
        atom = f"H 0 0 0; H 0 0 {R}"
        charge, spin = 1, 1
    mol = gto.M(atom=atom, basis=basis, charge=charge,
                spin=spin, verbose=0)
    mf = dft.UKS(mol)
    mf.xc = xc
    if tight:
        mf.conv_tol = 1e-10
        mf.max_cycle = 200
    mf.kernel()
    charges = mf.mulliken_pop(verbose=0)[1]
    return {
        "E": mf.e_tot,
        "converged": mf.converged,
        "cycles": mf.cycles,
        "charges": charges,
    }

def run_uhf_reference(system, R=None, basis="def2-tzvp", tight=True):
    # identical setup but scf.UHF; energy is exact for N=1
    ...
\end{verbatim}

\section{SCF stability analysis}

\begin{verbatim}
mf.kernel()
mf.stability()   # returns dict with internal/external stability info
\end{verbatim}

\todo{Record the exact output of \texttt{mf.stability()} for every
geometry in the master CSV, not just a boolean.}

\chapter{Raw data tables}

The master CSV is the authoritative source. Tables in
Chapter~\ref{ch:results} are reproductions of the corresponding CSV
slices. Any discrepancy between a table in this report and the CSV is a
reporting error to be fixed.

\todo{Generate this appendix automatically from the master CSV.}

% =====================================================================
\chapter*{Acknowledgements}
\addcontentsline{toc}{chapter}{Acknowledgements}
% =====================================================================

\todo{The AI coauthorship from the original drafts is removed in this
report. Add an acknowledgement line for any AI assistance received during
drafting, following the policy of the target journal.}

% =====================================================================
\begin{thebibliography}{99}
% =====================================================================

\bibitem{PZ1981}
J.~P.~Perdew and A.~Zunger,
\emph{Self-interaction correction to density-functional approximations
for many-electron systems},
Phys.\ Rev.\ B \textbf{23}, 5048 (1981).

\bibitem{PPLB1982}
J.~P.~Perdew, R.~G.~Parr, M.~Levy, and J.~L.~Balduz Jr.,
\emph{Density-functional theory for fractional particle number:
derivative discontinuities of the energy},
Phys.\ Rev.\ Lett.\ \textbf{49}, 1691 (1982).

\bibitem{LO1981}
E.~H.~Lieb and S.~Oxford,
\emph{Improved lower bound on the indirect Coulomb energy},
Int.\ J.\ Quantum Chem.\ \textbf{19}, 427 (1981).

\bibitem{PBE1996}
J.~P.~Perdew, K.~Burke, and M.~Ernzerhof,
\emph{Generalized gradient approximation made simple},
Phys.\ Rev.\ Lett.\ \textbf{77}, 3865 (1996).

\bibitem{SCAN2015}
J.~Sun, A.~Ruzsinszky, and J.~P.~Perdew,
\emph{Strongly constrained and appropriately normed semilocal density
functional},
Phys.\ Rev.\ Lett.\ \textbf{115}, 036402 (2015).

\bibitem{r2SCAN2020}
J.~W.~Furness, A.~D.~Kaplan, J.~Ning, J.~P.~Perdew, and J.~Sun,
\emph{Accurate and numerically efficient r$^2$SCAN
meta-generalized gradient approximation},
J.\ Phys.\ Chem.\ Lett.\ \textbf{11}, 8208 (2020).

\bibitem{B3LYP1994}
P.~J.~Stephens, F.~J.~Devlin, C.~F.~Chabalowski, and M.~J.~Frisch,
\emph{Ab initio calculation of vibrational absorption and circular
dichroism spectra using density functional force fields},
J.\ Phys.\ Chem.\ \textbf{98}, 11623 (1994).

\bibitem{Lowdin1963}
P.-O.~L\"owdin,
\emph{Symmetry dilemmas in Hartree--Fock theory},
Rev.\ Mod.\ Phys.\ \textbf{35}, 496 (1963).

\bibitem{CMSY2008}
A.~J.~Cohen, P.~Mori-S\'anchez, and W.~Yang,
\emph{Insights into current limitations of density functional theory},
Science \textbf{321}, 792 (2008).

\bibitem{CMSY2012}
A.~J.~Cohen, P.~Mori-S\'anchez, and W.~Yang,
\emph{Challenges for density functional theory},
Chem.\ Rev.\ \textbf{112}, 289 (2012).

\bibitem{CIDER2022}
A.~D.~Bystrom and B.~Kozinsky,
\emph{CIDER: an expressive, nonlocal feature set for machine learning
density functionals with exact constraints},
J.\ Chem.\ Theory Comput.\ \textbf{18}, 2180 (2022).

\bibitem{CIDER2024}
A.~D.~Bystrom and B.~Kozinsky,
\emph{Successor CIDER23X generation},
Phys.\ Rev.\ B (2024).

\bibitem{CIDER26}
A.~D.~Bystrom and B.~Kozinsky,
\emph{CIDER26: a machine-learned exchange--correlation functional for
molecules, solids, and surfaces},
arXiv:2608.21525 (2026).

\bibitem{Skala2025}
Microsoft Research AI for Science,
\emph{Skala: a machine-learned exchange--correlation functional}, (2025).

\bibitem{e3nn}
M.~Geiger and T.~Smidt,
\emph{e3nn: Euclidean Neural Networks},
arXiv:2207.09453 (2022).

\end{thebibliography}

\end{document}
